\section{Port Connector Taxonomy}
\label{sec:taxonomy}

\subsection{Definition and Scope}

A port connector segment has three defining characteristics that distinguish it from a standard T1 or T2 arterial:

\begin{enumerate}
  \item \textbf{Captive demand}: The truck volumes using the connector have no practical alternative. A shipper whose container arrives at the Port of Long Beach must use the I-710 (or the SR-47/SR-103 complex, which merge to the same interchange) to access the inland freight network. There is no routing alternative that avoids the constraint---unlike a T1 bottleneck at Donner Pass, where the I-40 alternate exists, however imperfect.

  \item \textbf{Drayage-dominated demand profile}: Port connector traffic is dominated by drayage trucks operating in 12-hour windows (typical port gate hours 6:00~am to 6:00~pm) rather than the 24-hour mixed-use traffic that characterizes T1 corridors. The peak-hour demand for a port connector can be 3--5 times the off-peak demand; for a comparable T1 segment, the peak-to-off-peak ratio is typically 1.5--2.0. Standard AADT-based capacity analysis, which spreads daily traffic across 24 hours, dramatically understates peak-hour congestion on port connectors.

  \item \textbf{Jurisdictional fragmentation}: A port connector typically traverses at least three jurisdictions---port authority land (which may have its own traffic management), city streets (controlled by the municipality), and state highway or interstate right-of-way (FHWA and state DOT). Investment in a port connector requires coordinated approval across all three, which creates transaction costs and delays that do not apply to pure interstate segments.
\end{enumerate}

\subsection{Why Port Connectors Are Invisible in Corridor Analysis}

Standard T1 tier analysis operates at the corridor scale: a corridor is defined as a continuous interstate segment, typically 50--300 miles, scored on dimensions including freight intensity, V/C ratio, climate exposure, and network centrality \citep{ROUTE_A1}. A 10-mile port connector segment is too short to be analyzed as an independent corridor; when the corridor-scoring algorithm encounters the I-710 segment, it averages the V/C across the full I-710 length (approximately 20 miles), diluting the 11-mile port connector's V/C 1.5 into a corridor-average of approximately 1.1---below the T1 priority threshold.

The ATRI Top Freight Bottlenecks report uses a different methodology: it ranks bottlenecks by total truck-delay-hours per year \citep{ATRI_costs2024}. A 10-mile segment generating 2,000 truck-hours of delay per day accumulates 730,000 truck-hours per year. An I-95 segment through New York generating 1,000 truck-hours per day over 50 miles accumulates 3,650,000 truck-hours per year and ranks higher. The ATRI methodology rewards length: a 5-times-more-congested short segment ranks below a moderately congested long segment. Port connectors lose this ranking systematically.

\subsection{Failure Mode Taxonomy}

We classify port connector failures into four modes that require distinct intervention types:

\textbf{Mode A -- Capacity Failure}: The connector has sufficient geometric standard (interstate or near-interstate) but insufficient lane capacity for the truck volumes generated by port growth. The I-710 Long Beach Freeway is the canonical case: it is a limited-access freeway with partial interchange control, but its 4--6 lane cross-section is insufficient for the current 35,000 trucks per day. The intervention is capacity expansion: managed lanes, widening, or parallel facility construction.

\textbf{Mode B -- Classification Failure}: The connector approach road does not meet interstate geometric standards even though it carries interstate-class truck volumes. SR-509 at the Port of Tacoma is the exemplar: a two-lane state route handling heavy port truck traffic with at-grade intersections, inadequate shoulder widths, and substandard interchange geometry. The intervention is geometric reconstruction to interstate standard---a more expensive and complex undertaking than pure capacity expansion.

\textbf{Mode C -- Resilience Failure}: The connector has a single chokepoint with no viable alternate. I-710 again qualifies: at its narrowest point (the Long Beach Gateway complex), there is no alternate routing available for the 35,000 trucks per day that use it. A bridge failure, major incident, or natural disaster on this segment shuts down the Port of Long Beach truck gate until the incident is cleared. The intervention is redundancy investment: a parallel connector or alternate interchange.

\textbf{Mode D -- Jurisdictional Fragmentation}: The connector's improvement is blocked by governance complexity rather than engineering or budget constraints. The Oakland port connector (I-880/Adeline Street approach) has been studied repeatedly but not improved because the improvement requires coordination between the Port of Oakland, the City of Oakland, Caltrans, and FHWA, each with different funding timelines and decision authorities. The intervention is jurisdictional: a designated port connector authority or a federal ``port connector'' funding category that pre-clears the approval pathway.

\subsection{The Top Ten Port Connectors: Baseline Data}

Table~\ref{tab:connectors} presents baseline data for the ten port connectors analyzed in this paper. Connector length is the distance from the port gate to the first T1 or T2 interstate interchange. Classification indicates the highest functional class of the connector segment.

\begin{table}[h]
\centering
\caption{Top 10 US Port Connectors: Baseline Data}
\label{tab:connectors}
\begin{tabular}{llccc}
\toprule
Port & Connector & Length & Classification & Failure Modes \\
\midrule
LA/Long Beach & I-710 (Long Beach Fwy) & 11 mi & T2 freeway & A, C \\
NY/NJ & I-278 / I-95 approach & 8 mi & T1 & A, D \\
Savannah GA & I-16/I-95 interchange & 6 mi & T2 & A \\
Houston Ship Channel & I-10/I-610 loop & 9 mi & T1 & A \\
Seattle/Tacoma & SR-509 / SR-518 & 7 mi & State route & B, C \\
Baltimore (Seagirt) & I-95/I-695 connector & 5 mi & T1 & A, D \\
Norfolk (VIT) & I-64 Hampton Roads & 12 mi & T1 & A, C \\
Charleston SC & I-526 Mark Clark Expwy & 9 mi & T2 & A, D \\
Miami (PortMiami) & I-395/MacArthur Cswy & 3 mi & T2 & A, D \\
Oakland & I-880/Adeline approach & 6 mi & T2 & A, D \\
\bottomrule
\end{tabular}
\end{table}

The connector lengths in Table~\ref{tab:connectors} are measured from port gate to first interstate on-ramp based on GIS analysis of the TIGER/Line road network \citep{Census_TIGER2023} overlaid with port gate locations from BTS Port Performance data \citep{BTS_airfreight2023}.
